\pdfbookmark[1]{Erläuterungen zu den Terminszeichnungen und den Beobachtungszahlen. 1837.}{erlaeuterungen-termine-1837}
\begin{center}\large\textit{Erläuterungen zu den Terminszeichnungen und den Beobachtungszahlen.}\par
[\textit{Im Auszug.}]\end{center}
\begin{center}\small Resultate aus den Beobachtungen des magnetischen Vereins.\quad 1837.\quad VIII.\end{center}

Im Jahr 1837 wurden sieben vierundzwanzigstündige Termine abgehalten, da zu den sechs gewöhnlichen noch ein außerordentlicher am 31. August hinzukam. Die in den ---\quad---\quad--- Tafeln mitgeteilten Zahlen enthalten 80 Beobachtungsreihen zu den Variationen der Deklination von 16 verschiedenen Beobachtungsorten, nämlich Altona, Augsburg, Berlin, Breda, Breslau, Kopenhagen, Dublin, Freiberg, Göttingen, Leipzig, Mailand, Marburg, München, Petersburg, Stockholm und Uppsala. Außerdem gingen bei uns noch einige andere Beobachtungsreihen ein, die zu spät eintrafen, um mit abgedruckt werden zu können.
\begin{center}\small
---\enspace---\enspace---\enspace---\enspace---\enspace---\enspace---\enspace---\enspace---\enspace---\enspace---\enspace---\enspace---\enspace---\enspace---
\end{center}
Den Anlass für einen außerordentlichen Termin am 31. August gab die Nachricht, dass Herr Professor \textsc{Parrot} an diesem Tag (wie auch an einigen vorhergehenden) die Variation der magnetischen Deklination am Nordkap beobachten würde. Die Einladung zur Teilnahme wurde daher verbreitet, soweit es die Kürze der Zeit zuließ. Dadurch gingen recht interessante Beobachtungsreihen ein, doch die Beobachtungen vom Nordkap selbst sind uns bisher nicht bekannt geworden. ---\quad---\quad---\quad---\quad---\quad---\quad---

\Needspace{23\baselineskip}
Für den Novembertermin wurde die sonst befolgte Regel geändert und der Termin auf den 13. verlegt. Dies geschah nach einem Gespräch mit Herrn von \textsc{Humboldt} über die Möglichkeit, dass an den Tagen des Monats, die in mehreren früheren Jahren durch eine außerordentliche Zahl von Sternschnuppen aufgefallen waren,
\SourceLeaf{583}vielleicht auch ungewöhnliche magnetische Bewegungen auftreten könnten. Diese Erwartung hat sich jedoch insofern nicht bestätigt, als die magnetischen Bewegungen während dieser vierundzwanzig Stunden zwar sehr beträchtlich, aber doch nicht größer waren als bei vielen früheren Terminen zu jeder anderen Jahreszeit. Dagegen wurden am vorhergehenden und am folgenden Abend an mehreren Orten sehr starke und schnell wechselnde Anomalien der magnetischen Deklination beobachtet. Zwischen ihnen und den Sternschnuppenerscheinungen darf man jedoch keinen Zusammenhang annehmen, da jene Anomalien nur die gewöhnlichen Begleiter von Nordlichtern sind und in diesen beiden Nächten tatsächlich sehr leuchtende Nordlichter auftraten\expandafter\def\csname HistoricalNoteText583\endcsname{Die magnetischen Beobachtungen vom 12. November wurden uns aus Uppsala, Leipzig, Breslau und Mailand mitgeteilt, die vom 14. November aus Uppsala, Dublin, Berlin, Breslau und Mailand. Die Ähnlichkeit der Bewegungen ist hier an einigen Stellen unverkennbar, an anderen nur schwach zu erkennen. Doch hier lässt sich erneut feststellen, was schon bei anderer Gelegenheit bemerkt wurde. Unter solchen Umständen wechseln die Bewegungen viel zu schnell, als dass Beobachtungen im Abstand von fünf Minuten oder gar in noch größeren Abständen ein getreues Bild davon geben könnten.}\HistoricalNote{583}.
\begin{center}\small
---\enspace---\enspace---\enspace---\enspace---\enspace---\enspace---\enspace---\enspace---\enspace---\enspace---\enspace---\enspace---\enspace---\enspace---\enspace---
\end{center}
Bei den Terminen im Juli, August und November wurden in Göttingen nun auch die Schwankungen der Intensität mit dem Bifilarmagnetometer vollständig beobachtet. In die Tafeln wurden jedoch nicht die unmittelbar beobachteten Skalenwerte selbst aufgenommen. sondern ihre Differenzen zum größten Wert, der bei dem jeweiligen Termin vorkam. Da bei den beiden ersten Terminen diejenige Querlage vorlag, in der zunehmenden Skalenwerten abnehmende Intensitäten entsprechen, geben die Zahlen hier an, um wie viel die jeweilige Intensität größer war als die kleinste des Termins, gemessen in Einheiten, von denen beim Julitermin 22000 auf die kleinste Intensität selbst kommen. Solange derartige Beobachtungen nur an einem einzigen Ort gemacht werden, kommt es vorerst nicht darauf an, den \textit{absoluten} Wert der Skalenteile so genau wie möglich anzugeben. Für den Augusttermin wurden daher zu diesem Zweck keine neuen Bestimmungen vorgenommen. Vor dem Novembertermin war dies dagegen geschehen. Die geänderte absolute Zahl hängt mit dem Verlust zusammen, den der Magnetismus des Stabs in den vier Monaten erlitten hatte. Zu bemerken ist nur, dass die Zahlen beim Novembertermin angeben,
\SourceLeaf{584}um wie viel die jeweilige Intensität kleiner ist als die größte, wobei diese selbst als $=18290$ angenommen wird. Bei diesem Termin hatte der Stab nämlich die entgegengesetzte Querlage, bei der Intensität und Skalenwerte zugleich zunehmen. Damit die Zahlen dieselbe Bedeutung hätten, hätte man die einzelnen unmittelbar beobachteten Skalenwerte daher mit dem Minimum statt mit dem Maximum vergleichen müssen. Das wurde versehentlich übersehen.

Auf den Tafeln ---\quad---\quad---\quad--- sind die beobachteten Intensitätsänderungen grafisch dargestellt. Zunächst in einer Kurve, unter der die Kurve für die gleichzeitigen Deklinationsänderungen in Göttingen wiederholt ist. Dadurch wird die obige Bemerkung [S.\ 365 d.\ B.] augenfällig, dass zur Zeit starker Störungen der Deklination meist auch starke Anomalien der Intensität auftreten. Zweitens ist auch der Verlauf der Änderungen beider Elemente in einer Kurve zusammengefasst. So erhält man ein anschauliches Bild der Änderungen der horizontalen Komponente der erdmagnetischen Kraft während jedes Termins. Um Verwirrung durch die vielen Kreuzungen zu vermeiden, mussten die Bewegungen beim Juli- und Augusttermin allerdings in zwei Stücken, beim Novembertermin in drei Stücken gezeichnet werden. Der Übersichtlichkeit halber ist außerdem jedes Stück zur einen Hälfte mit durchgezogenen Linien, zur anderen punktiert dargestellt. Nach dem, was bereits oben [S.\ 366 d.\ B.] bemerkt wurde, werden diese Darstellungen keiner weiteren Erläuterung bedürfen.

Zu den Ergebnissen selbst können hier nur noch einige Bemerkungen Platz finden. Die außerordentlich große Ähnlichkeit der gleichzeitigen Deklinationsbewegungen an verschiedenen Orten, meist bis in die feinsten Abstufungen, bestätigt sich hier wieder ebenso schön wie bei den Beobachtungen des vorhergehenden Jahres. Doch werden hin und wieder auch schon erhebliche Unterschiede erkennbar, besonders bei den Terminen, an denen die Beobachtungen ein noch größeres Gebiet abdecken. Dieses Gebiet scheint aber noch immer zu klein und die Zahl weit auseinanderliegender Orte zu gering, um darauf bereits Schlüsse über den Sitz der Ursachen der einzelnen Bewegungen zu stützen. Immerhin könnte eine nähere Betrachtung mancher einzelner Bewegungen, zumal bei den Terminen, an denen in Göttingen zugleich die Intensitätsänderungen beobachtet wurden,
\SourceLeaf{585}Anlass zu allerlei Bemerkungen und sogar zu allgemeinen Betrachtungen geben. Darin wollen wir unseren Lesern jedoch nicht vorgreifen. Stattdessen möchten wir daran erinnern, dass man bei allen erkennbaren Unähnlichkeiten vor allem die äußeren Umstände sorgfältig abwägen muss, ehe man sie zur Grundlage gewagter Vermutungen macht. Als Beispiel kann die kleine Erhöhung dienen, die man in den grafischen Darstellungen des Augusttermins ---\quad---\quad---\quad--- um $18^{\mathrm u}5'$ an den meisten Beobachtungsorten, am stärksten am nördlichsten, Uppsala, bemerkt. Dass sie in Dublin fehlt oder nur eine schwache Spur davon sichtbar ist, ist allerdings bemerkenswert, da kein Grund besteht, die Richtigkeit der Beobachtung selbst zu bezweifeln. Das würde uns, zumal im Zusammenhang mit der vollständigen Erscheinung in Göttingen ---\quad---, Anlass zu interessanten Betrachtungen geben, wenn es überhaupt angemessen wäre, uns hier schon darauf einzulassen. Dass diese Erhöhung aber auch in Kopenhagen fehlt, ist völlig bedeutungslos, weil in Kopenhagen bei diesem Termin nur von 10 zu 10 Minuten, also um $18^{\mathrm u}5'$ gar nicht beobachtet wurde.

Zu den labyrinthischen Formen, die uns die magnetischen Beobachtungen zeigen, wenn Deklinations- und Intensitätsbewegungen in einer Kurve ---\quad---\quad---\quad---\quad--- vereinigt werden, enthalten wir uns hier jeder Bemerkung, allein weil begründete Hoffnung besteht, dass bald viel reicheres Material zur Verfügung stehen wird. Wer sich inzwischen selbst schon an Betrachtungen darüber versuchen möchte, braucht sich wenigstens nicht durch Zweifel an der Wirklichkeit der vom Bifilarmagnetometer angezeigten Intensitätsbewegungen davon abhalten zu lassen. In der Tat sind solche Zweifel ganz unstatthaft geworden, nachdem bereits beim Märztermin des gegenwärtigen Jahres 1838 außer in Göttingen auch an \textit{drei anderen Orten} die gleichzeitigen Intensitätsbewegungen mit ähnlichen Bifilarapparaten beobachtet wurden und eine ebenso bewundernswerte Übereinstimmung zeigten, wie wir sie seit vier Jahren bei den Deklinationsbewegungen zu finden gewohnt sind. Näheres dazu wird jedoch der Bekanntmachung der Ergebnisse der Beobachtungen von 1838 vorbehalten bleiben müssen.

